\CSBHeading{Anschluss in Band 11: Gleichmächtigkeit}
Band~11 führt die Schreibweise \(A\CardLeq B\) für die Existenz
einer Injektion von \(A\) nach \(B\) ein
(\FormulaRefAuto{CardLeqDef}[def]). Die Gleichmächtigkeit
\(A\EqCard B\) bedeutet die Existenz einer Bijektion
(\FormulaRefAuto{Gleichmächtigkeit}[def]).

Gelten \(A\CardLeq B\) und \(B\CardLeq A\), so gibt es passende
Injektionen \(F\colon A\inj B\) und \(G\colon B\inj A\).
Der eben bewiesene Satz liefert eine Bijektion \(H\colon A\bij B\)
und damit \(A\EqCard B\). Umgekehrt liefern eine Bijektion und ihre
Umkehrfunktion die beiden Injektionen. Somit gilt
\[
  A\EqCard B\quad\Longleftrightarrow\quad
  A\CardLeq B\land B\CardLeq A.
\]
Die formale Übertragung steht in Band~11 unter
\FormulaRefAuto{CantorSchroederBernstein}[thm] und
\FormulaRefAuto{EqCardMutualCardLeqEquiv}[thm]. Dafür wird die
Konstruktion von \(C\) nicht erneut benötigt. Dieser Anschluss ist
eine spätere Anwendung; der Beweis in Band~08 setzt diese Begriffe nicht voraus.
